Durante el desarrollo del proyecto de residencia profesional se aplican diversas competencias profesionales, destacando aquellas relacionadas con la ingeniería de datos. En primer lugar, se demuestra capacidad técnica para el diseño y creación de pipelines de datos automatizados, garantizando una correcta integración y procesamiento de la información. Asimismo, se desarrolla la habilidad para analizar requerimientos y traducirlos mediante el modelado de datos, estructurando esquemas analíticos que resuelven problemas específicos de la organización.

Por otro lado, se fortalece el trabajo en equipo y la comunicación efectiva con diferentes áreas de la empresa, asegurando la alineación de los objetivos técnicos con las necesidades operativas. Finalmente, se aplica el pensamiento crítico para la toma de decisiones en el diseño de las transformaciones de datos y el aseguramiento de la calidad de la información.
